\documentclass[11pt]{article}
\usepackage[margin=1in]{geometry}
\usepackage{parskip}

\begin{document}

Most research organizations either publish papers or ship products. Basis does both, and treats open-source software as the primary research medium rather than a side artifact. That's rare, and it's why I'm interested.

My background spans robotics research (PhD, Georgia Tech) and building production systems as a founder. I've found the best work happens when foundational thinking meets real-world constraints. At Pytheia and Roostr, that meant turning theory into systems that actually ran in production with customers, latency, reliability, and integration constraints. Basis seems to operate there by design.

The R-ADA project is particularly compelling. My dissertation was on optimization for spacecraft control systems: designing controllers that satisfy hard constraints under uncertainty. R-ADA's co-design of morphology and control, using probabilistic programming to bridge sim-to-real, feels like a natural extension of that line of thinking. I'd be excited to contribute.

More broadly, I'm drawn to Basis's bet that general principles of reasoning, from programming languages, compilers, and probabilistic inference, can unlock problems that brute-force scaling cannot. That's the kind of foundational work I want to be doing.

\vspace{1em}
Mark Mote

\end{document}
